%!TEX root = main.tex
% ==============================================================================
% İSTANBUL ÜNİVERSİTESİ SOSYAL BİLİMLER ENSTİTÜSÜ
% GİRİŞ BÖLÜMÜ (Madde 16/a, b ve Madde 17)
% Giriş, metin kısmının ilk öğesidir; Arap rakamlarıyla numaralandırma buradan başlar.
% Giriş bölümü numaralanmaz (Madde 17).
% ==============================================================================
\chapter*{GİRİŞ}
\addcontentsline{toc}{chapter}{GİRİŞ}

Giriş kısmında, önsözde belirtilen hususlar tekrar edilmemek üzere, çalışmada çözümlenmesi amaçlanan bilimsel sorun/sorunsal etraflıca tanımlanır. Bu doğrultuda araştırmada yararlanılan kavramsal çerçeve, kuramsal dayanaklar, yöntem, teknik ve varsa kullanılan paradigmalar ayrıntılı bir biçimde açıklanır (Madde 16/b).

Tez metninde satırlar iki yana yaslanmış ve 1,5 satır aralığı ile yazılmıştır (Madde 4/b). Paragraflar arasında gereksiz dikey boşluk bırakılmamış, standart 1 cm paragraf girintisi uygulanmıştır. 

Araştırma konusuna ilişkin genel değerlendirmeler ve tezin kurgusu Giriş kısmında ortaya konulduktan sonra, tezi sonuca ulaştıracak temel bölümlere geçilmektedir.

